\documentclass[12pt,english]{article}
\usepackage[utf8]{inputenc}
\usepackage{tgpagella} %
\usepackage{mathpazo}  %
\usepackage[left=1.5in,right=1.5in,top=1.5in,bottom=1.5in]{geometry}
\usepackage[round,sort&compress]{natbib} %
\usepackage{url} %
\usepackage[hyperpageref]{backref} %
\usepackage[pagebackref]{hyperref} %
\hypersetup{
  colorlinks=true, %
  urlcolor=purple,    %
  linkcolor=blue,  %
  citecolor=violet,   %
}
\usepackage{bm}
\usepackage{indentfirst}
\setcitestyle{aysep={}} 
\usepackage{amsmath}
\usepackage{mathtools}
\usepackage{mathrsfs}
\usepackage{tcolorbox}
\usepackage{amssymb}
\usepackage{eurosym}
\usepackage{amsfonts}
\usepackage{enumerate}
\usepackage{babel}
\usepackage{graphicx}
\usepackage{caption}
\usepackage{supertabular}
\usepackage{tabularx}
\usepackage{float}
\usepackage{dsfont}
\usepackage{fancyvrb}
\usepackage{verbatim}
\usepackage{enumitem}
\usepackage{setspace}
\usepackage{comment}
\usepackage{subcaption}
\usepackage{tikz}
\usepackage{gensymb}
\usepackage{textcomp}
\usepackage{placeins} %
\usepackage{tabulary}
\usepackage{tabularx}
\usepackage{booktabs}
\usepackage{fullpage}
\usepackage{morefloats}
\usepackage{makecell}
\usepackage{lscape}
\usepackage{pdflscape}
\usepackage{longtable}
\usepackage{rotating}
\usepackage{fancyhdr}
\usepackage{tocbibind} %
\usepackage{titletoc} %
\usepackage[export]{adjustbox}
\usepackage[anythingbreaks]{breakurl} %
\usepackage{multicol}
\usepackage{lineno}
\usepackage{endnotes}
\usepackage{ragged2e}
\newcolumntype{P}[1]{>{\RaggedRight\hspace{0pt}}p{#1}} %
\linenumbers
\makeatletter
\renewcommand\tableofcontents{%
    \@starttoc{toc}%
}
\makeatother
\newsavebox\ltmcbox %
\renewcommand{\floatpagefraction}{.99}
\newenvironment{stretchpars}{\par\setlength{\parfillskip}{0pt}}{\par} %

\newcommand{\bo}[1]{\textbf{#1}}
\newcommand{\gras}[1]{\textbf{#1}}

\title{ITMO Regulations to Ratchet up Ambition
} %
\date{\today} %

\begin{document}

\sloppy
\maketitle

\begin{abstract} %
The %
international climate policy regime is insufficient to achieve the Paris Agreement temperature target. Although Internationally Transferred Mitigation Outcomes (ITMOs) aim to make mitigation more efficient, their unfettered use could weaken the domestic ambition of countries buying them and insufficiently compensate seller countries. To bolster ambition and align Nationally Determined Contributions (NDCs) with the Paris target, this paper proposes that a coalition of willing countries commit to additional rules governing ITMO use and NDC definition. Basic rules include the harmonization of NDC targets using a comparable metric and requiring that no ITMO be purchased from a country missing its NDC. Then, additional rules with increasing degrees of ambition are described. First, restrictions based on national assessments are discussed, in which coalition members commit not to buy ITMOs from countries with unambitious NDCs. Next, it is argued that NDCs should be assessed for the coalition as a whole rather than country-by-country. Participating countries would cooperate in determining their NDCs to ensure joint alignment with the Paris target and commit not to exchange ITMOs with countries outside the coalition. 
The article concludes that combining a tightly regulated ITMO coalition with ambitiously expanded Just Energy Transition Partnerships offers a viable pathway for effective international decarbonization.
\end{abstract}

\textbf{Keywords}: Climate policy; carbon price; ITMOs; carbon trading.%

\textbf{JEL}: Q56; F38; H23; Q54; H87; F64; Q58; F53; F35.

\clearpage

\section*{Introduction}
In 2015, governments universally adopted a global temperature target at the Paris Agreement. While countries have developed various commitments, mechanisms, and initiatives to stabilize the climate, they are still insufficient to achieve the target. This paper proposes new rules for a coalition of countries to institutionalize higher ambition within cross-border carbon trading.

At COP29 in 2024, the international community finalized the rules for operationalizing Article 6.2 of the Paris Agreement, which allows countries to trade Internationally Transferred Mitigation Outcomes (ITMOs). While Article 6 aims to foster higher ambition by directing mitigation funding where it is most effective, the unfettered use of ITMOs risks weakening climate action: developing countries may set unambitious NDCs to sell more ITMOs, while developed countries may buy ITMOs as a cheaper substitute for domestic action. To the extent that NDCs remain collectively inconsistent with the Paris target, ITMOs will be underpriced (unless their use is restricted). While \cite{michaelowa_additionality_2019} proposed restrictions on ITMO use, I propose more stringent rules to strengthen ambition.

I start by describing existing proposals by \cite{michaelowa_additionality_2019} and \cite{la_hoz_theuer_when_2019}, then propose different options for leveraging ITMOs to enhance climate ambition, with increasing levels of stringency (summarized in Table \ref{tab:proposals}). Among the rules I propose, some rely on national assessments of NDCs: principled buyers would commit not to buy ITMOs from a country with an unambitious NDC. Besides, the conditional NDC target of a Global South country would supersede its unconditional target in adequacy assessments once the country's conditions in terms of climate finance are met. Hence, when a country sets its unconditional target above its projected emissions and its conditional target below them, buyers would need to provide climate finance to that country before purchasing ITMOs from it, as the seller's NDC would otherwise not be ambitious enough. Finally, I detail the limitations of relying on national assessments and make the case for a coalition of the willing to jointly align its NDCs with the Paris target and commit not to exchange ITMOs with countries outside the coalition. %
\section{Existing proposals\label{subsec:existing}}
\cite{michaelowa_additionality_2019} propose verifying that ITMOs are additional, meaning they correspond to emission reductions relative to a counterfactual scenario without such ITMOs. The authors suggest a two-step procedure wherein a new ``Article 6 Supervisory Board'' would first check whether the seller's NDC is more ambitious than the (frequently updated) business-as-usual (BAU) trend. %
If so, the second step would be unnecessary. Otherwise, the seller's unambitious NDC generates ``hot air'', requiring ITMO to be tested at the project level. In this second step, the additionality of each specific activity financed by the ITMO would be assessed individually. 

\cite{la_hoz_theuer_when_2019} discuss restricting ITMOs to emission reductions below the BAU. They compute BAU emissions for a range of countries using four potential definitions for the BAU: extending past emissions levels, extending emission intensity, or projecting their respective historical trend. They show that none of these definitions is satisfactory, as each would still generate \textit{hot air} for some countries where the BAU exceeds emissions as projected by Climate Action Tracker. The authors propose instead to limit the quantity of ITMOs that a country can sell to a fixed share of its emissions (e.g. 1\% or 5\%). They acknowledge that quantity limits would still allow some hot air but argue that such limits can be calibrated to strike a balance between reducing hot air and exploiting the gains from %
trade. %
While \cite{la_hoz_theuer_when_2019} show that simple definitions of the BAU are ill-suited for regulating ITMOs, \cite{michaelowa_additionality_2019} propose delegating the determination of the BAU to independent experts so that NDC ambition can be accurately assessed.\footnote{Expert projections of BAU emissions already exist, see e.g. \cite{den_elzen_are_2019,den_elzen_impact_2023,den_elzen_infographics_2024,unep_emissions_2025}.} In both papers, the authors agree that an NDC can be assessed at the national level, by comparing it to the country's projected emissions, and that ITMOs should be allowed whenever the NDC is more ambitious than properly projected emissions. 

\section{Desirable paths to regulate carbon trading\label{subsec:desirable}}

Principled %
buyers of ITMOs could also employ other national criteria to assess the adequacy of a seller's NDC. Here are different original proposals, ordered by degree of stringency.

\subsection{A necessary precondition: harmonized NDCs} %
Currently, few constraints apply to the definition of NDCs. Countries independently choose their sectoral and gas coverage, the methods for converting different gases into CO$_\text{2}$-equivalent, whether to use an absolute or carbon intensity target, etc. As a result, NDCs are not harmonized. Researchers must make assumptions (e.g. on GDP growth or LULUCF emissions) to express NDCs using a comparable metric, leading to significant differences between estimates of NDC targets across research teams.\footnote{Estimates of emissions implied by NDC targets are provided by \href{https://climateactiontracker.org/}{Climate Action Tracker}, \href{https://www.climateandforests-undp.org/plant}{UNDP}, \href{https://1p5ndc-pathways.climateanalytics.org/}{Climate Analytics}, \href{https://www.climate-resource.com/tools/ndcs}{Climate Resource}, \href{https://zenodo.org/records/10875617}{PBL}, \href{https://www.climatewatchdata.org/}{Climate Watch}, \cite{nascimento_comparing_2024,den_elzen_updated_2022,unep_emissions_2025}.} %

A necessary precondition to assessing NDCs is the strengthening of reporting requirement. NDCs should cover all Kyoto gases and be expressed in absolute terms using harmonized conversion factors between gases (namely GWP100 or GWP500, the most up-to-date being from the IPCC AR6). Furthermore, they should encompass all sectors---broken down individually, or at least separating LULUCF from non-LULUCF emissions. Ideally, NDCs should also specify a country's future cumulative emissions, its planned emission trajectory, and the intended use of ITMOs. These enhanced reporting requirements would allow for the proper assessment of NDCs and enable conditioning the use of ITMOs on the achievement and adequacy of NDCs. 

\subsection{A sine qua non condition: no ITMOs from failed NDCs} %
As a minimal requirement, countries should be allowed to sell ITMOs only up to the amount by which their emissions fall below their NDC targets. Consequently, buyers should refrain from purchasing ITMOs from a country that is failing to meet its NDC target. %
This principle should also apply to countries that have weakened their emission targets (instead of ratcheting up ambition), resulting in emissions higher than planned under a previous target (\href{https://climateactiontracker.org/climate-target-update-tracker/list-non-updating-countries/}{Climate Action Tracker} gives examples of such countries).

\subsection{Principled buyers: no ITMOs from unambitious NDCs} %
Furthermore, to strengthen the additionality requirement proposed by \cite{michaelowa_additionality_2019}, %
principled buyers should refuse to buy ITMOs from a country whose NDC target is above the BAU\footnote{The BAU is understood as a projection of emissions which accounts for the most recent technological developments.} 
  %
or evidently incompatible with the Paris target, even when it is below the BAU.
Such incompatibility could be safely assumed when a country's NDC target exceeds both its cost-optimal emissions required to limit global warming to 2\textdegree{}C and its equal per capita share of global emissions under that 2\textdegree{}C scenario. %
Indeed, no credible burden-sharing rule could then justify the NDC's adequacy with the Paris temperature target %
 \citep{van_den_berg_implications_2020,robiou_du_pont_equitable_2017,hohne_regional_2014}. %
Given that there are multiple ways to define a cost-optimal 2\textdegree{}C scenario (as the carbon budget depends on the probability of achieving the 2\textdegree{}C target, while the trajectory depends on the model used), simple substitutes could be employed instead. In particular, buyers could refuse to acquire ITMOs from countries with per capita emissions %
above the world average. %

The first exchange of ITMO that has been finalized was between Thailand and Switzerland. %
Thailand appears to meet all aforementioned criteria, suggesting that Switzerland could be considered a principled buyer. However, until the sum of all countries' NDCs aligns with the Paris target, the price of ITMOs might be too low. In other words, the above requirements are necessary but not sufficient conditions to close the ambition gap, understood here as the gap between the Paris target and the sum of NDCs \citep{unep_emissions_2025}. %
\subsection{Adequate climate finance unlocking conditional NDCs} %
While these requirements would reduce the amount of hot air, they would exclude from the trade of ITMOs low-emitting countries that set their NDC target above the BAU,\footnote{For example, \cite{van_den_berg_implications_2020} and \cite{unep_emissions_2025} conclude that India has set its 2030 NDC above the BAU (note that the BAU trajectory was computed by the authors; India's NDC does not specify a BAU). %
} 
at (what they consider to be) their fair share. These potential sellers of ITMOs could join forces and use their market power to negotiate guarantees of grant-based climate finance at scale. In particular, these countries could commit to setting their NDC target below the BAU (and selling ITMOs) in exchange for a commitment from buyers to fairly allocate taxing rights of new international levies, or to finance debt relief or Just Energy Transition Partnerships (JETPs). %

There is another way sellers would incentivize principled buyers to provide sufficient climate finance. Until now, I conflated conditional and unconditional NDC targets. 
Yet, Global South countries often include both types of target: the conditional target sets more stringent emission reductions than the unconditional one, but applies only at the condition that the country receives a specified amount of climate finance from developed countries. Let me call \textit{unidirectional climate finance} all climate finance flows that do not rely on ITMOs. When assessing a country's NDC (to determine whether ITMOs can be purchased from that country), the conditional NDC target should replace the unconditional target in the adequacy assessment %
once the country's conditions (in terms of climate finance) are met, counting only unidirectional climate finance. Therefore, if a country sets its unconditional target above the BAU and its conditional target below the BAU, principled buyers would have to provide climate finance to that country in order to purchase ITMOs from it, otherwise the seller's NDC would not be ambitious enough.

\section{Limitations of the previous options\label{subsec:limitations}}

\subsection{Issues with national assessments of NDCs}
The aforementioned proposals %
rely on assessing NDCs against national BAUs. However, this approach raises two issues. 
First, even if a country's NDC target is set below its BAU, to the extent that NDCs do not align with the Paris temperature target when aggregated at the global level, the ambition gap will persist. %
Consequently, both the demand for ITMOs and their price will remain too low. Second, given that %
grant-based climate finance falls short of their demands, Global South countries could legitimately set up unconditional NDC emission targets above their BAU emissions to use extra emission space to sell ITMOs. Actually, insofar as NDC targets represent how countries should share the burden of emission reductions, low-income countries already claim less than their fair share of emissions.\footnote{Based on my own estimates, the 2030 NDC targets of Global South countries (excluding China) total 13.5~GtCO$_\text{2}$, that is 2.8~tCO$_\text{2}$ per capita, which is below equal per capita emissions aligned with the Paris target.} %
Therefore, the ratcheting up of ambition should arguably be borne by industrialized countries (including China), which should decarbonize faster and provide climate finance abroad. For these reasons, assessing an NDC against the national (or project-based) BAU is neither a fair nor effective way to close the ambition gap. %

A rule restricting the sale of ITMOs to cases where the NDC target is below the BAU would be satisfactory only at a global scale or within a large coalition, since this is the scale at which there is consensus that the BAU is inadequate.
Conversely, at the national level, some countries may have already implemented stringent climate policies %
such that their BAU emissions are ambitious enough, while others would legitimately set an NDC target above their BAU since they deserve to sell ITMOs. In both cases, a national comparison of the NDC to the BAU would be inappropriate. 
The key problem with ITMOs is that the \textit{global} emissions implied by current NDCs (or current trends) do not align with the Paris temperature target \citep{unep_emissions_2025,rogelj_credibility_2023}. 
Meanwhile, given the disagreement over what constitutes a fair share of the global carbon budget, countries cannot agree on whether a given NDC complies with the Paris target or %
claims an excessive carbon budget. However, if a critical mass of countries agreed on a common norm, such as a burden-sharing principle or on a decision rule to assess whether the global carbon budget is fairly shared, then the two issues of effectiveness %
and fairness could be resolved. 

\subsection{Limitations of restrictions to the use of ITMOs}
As \cite{mehling_governing_2019} explain, restricting the use of ITMOs involves a trade-off between limiting hot air and exploiting gains from trade. Furthermore, the restrictions discussed above present additional limitations. They rely on arbitrary cutoffs, which create undesirable discontinuities in the ability to sell ITMOs. Additionally, they do not provide an institution for states to negotiate %
common ground regarding overall ambition or the scale of climate finance. 

A coalition of willing countries could %
establish precise norms for NDC adequacy and climate finance contributions. The adoption of such norms by a large group of countries would alleviate the need for restrictions on ITMO exchanges between them (beyond the sine qua non condition) and overcome the aforementioned limitations.

\section{A coalition of the willing to jointly define NDCs\label{sec:itmo}}

\subsection{The case for a joint definition of NDCs}\label{subsec:itmo_proposal}

A coalition of the willing could even go further and submit a joint NDC %
(just as the European Union does). 
The common norm would be used to verify that the joint NDC complies with the global target and thereby closes the ambition gap. %
This coalition could also allocate its aggregate NDC between countries in a way deemed acceptable and fair, potentially allowing certain countries to get a target higher than their current policies emissions.

Ideally, the joint NDC would include a carbon budget aligned with the Paris target, broken down into yearly national targets. That way, the adequacy of emission trajectories could be verified jointly and country-by-country, 
and the system could easily evolve into a fully-fledged compliance carbon market. 
Alternatively, the joint NDC would define a minimum emission reduction rate, aligned with the Paris target. Initially, this rate could be expressed in terms of emission intensity, consistent with the practice of emerging economies. For example, the coalition's GHG emissions relative to output or final energy use\footnote{Although emission intensity is usually defined relative to GDP, it may be more advisable to define it relative to final energy use, as emissions may decouple from GDP more easily than from energy use, and GDP indicators may be more susceptible to manipulation.} %
would need to decrease by 2\% each year. In the medium term, the reduction should be defined in absolute terms.
The countries most likely to join such a coalition are those with moderate emissions. There is therefore a tension between setting the coalition's carbon budget based on a cost-optimal allocation (favoring large emitters outside the coalition) or an egalitarian per capita allocation (which might not sufficiently strengthen decarbonization ambition). Negotiations within the coalition would be essential to defining fair shares.
\subsection{A coalition for the use of ITMOs\label{subsec:itmo}}

Here is how negotiations could shape an agreement along the lines sketched above. A coalition of countries could agree to jointly define their NDCs and exchange ITMOs exclusively among themselves to ensure that their carbon trading does not undermine the ambition of the Paris Agreement. 
Once a large coalition of countries agree on the broad vision, this coalition could be taken as given. The coalition's emissions targets would be set below its projected emissions, with a gradually increasing wedge between current policies emissions and the target. Specifically, the 2026 target would correspond to the coalition's emissions, %
while the 2027 target would be slightly below the current policies projection, increasing the divergence in 2028 and beyond. This would realistically scale up the additional decarbonization effort over time 
until it reaches an emission reduction rate aligned with the Paris target. 
The resulting emission trajectory of the coalition should not exceed its equal per capita share of a global carbon budget achieving 2\textdegree{}C with a 67\% probability,\footnote{That is equivalent to 1.9\textdegree{}C with a 50\% probability \citep{forster_indicators_2025}.} and would ideally be lower. %

After determining the coalition's trajectory, the disaggregation into national targets would be negotiated. It would be useful to allocate national targets starting from a focal point, such as  %
the minimum between the country's current policies projected emissions, its NDC targets, and an equal per capita share of the coalition's emission trajectory. 
The proposed focal point can be justified by the principle that, to strengthen ambition, a country's target should not be higher than projected, committed, and \textit{fair-share} emissions at once, and by the norm that the equal right to emit per capita is the fairness benchmark. %

Given that the equal per capita share is above projected emissions %
and the NDC for some countries, the sum of focal point emissions will be lower than the coalition's emission target, and some extra emission rights will be available. To the extent that projected (and committed) emissions are far lower than the fair share in some low-income countries, and that emission reductions below the NDC are very costly for some middle-income countries, the extra rights should be allocated to these low-income or fossil-dependent countries. While the precise allocation rules need to be negotiated, I propose a specific allocation into yearly national targets in the working paper version [anonymized for review]. %

If new countries join the coalition, they could propose a new allocation of emission targets, adopted only if accepted by a majority within the coalition and provided it does not lead to reduced emission coverage (resulting from unsatisfied countries leaving the coalition) or a lower projected ITMO price (which would indicate reduced ambition). %

A scientific council would assist the coalition by determining current policies emissions, modelling climate and economic effects, and proposing target allocations. Each participating country would designate a team of scientists (which could overlap), with voting rights proportional to the country's population in the event of disagreement.

\subsection{Feasibility of the previous proposal\label{subsec:feasibility}}

While the last proposal is the most ambitious analyzed here, I argue that it constitutes the most feasible pathway to guarantee emission reductions in line with the Paris target within a large coalition of countries. 

First, relying on a coalition parallel to the UNFCCC would enable the achievement of ambitious agreements. %
Second, %
academic surveys show that majorities around the world would genuinely accept international climate policies, even knowing the cost to themselves \citep{fabre_majority_2025,fabre_public_2025,andre_globally_2024}. %
Third, the proposals presented in this article are grounded in the Paris Agreement (and in particular Article 6.2) and they would operationalize the long-standing CBDR principle. These proposals respond to a renewed interest in well-regulated international carbon markets, as evidenced by two initiatives launched at COP30: the \textit{Open Coalition on Compliance Carbon Markets}, which already includes China, the EU, Brazil, and fifteen other countries; and the \textit{Article 6 Ambition Alliance}, a set of ten countries including Switzerland, Norway, and Germany, which intend to finance or host some emission reduction activities without use of ITMOs towards their NDCs.

\begin{table}
  \caption{Proposals to strengthen ambition using ITMOs or NDCs. Restrictions based on national assessments of NDCs are denoted with \textit{Nat}, those based on BAU trajectory are denoted with \textit{BAU}. \label{tab:proposals}}
  \makebox[\textwidth][c]{%
\begin{tabular}[t]{P{5.9cm}P{.9cm}P{13.2cm}}
  \toprule
  Policy & Type & Description \\
  \midrule 
  \textbf{Harmonized NDCs} &  & NDCs should cover all sectors and GHGs, use an absolute target, and specify cumulative emissions. \\
  \textbf{ITMOs should be additional} \citep{michaelowa_additionality_2019} & Nat, BAU & ITMOs should either be sold by a country whose NDC target is below the BAU, or their additionality verified at the project level. \\
  \textbf{ITMOs should be limited} \citep{la_hoz_theuer_when_2019} & Nat & ITMOs should not exceed a fixed share (e.g. 1\% or 5\%) of the seller's emissions. \\
  \textbf{No ITMOs from failed NDCs} & Nat & ITMOs sold should not exceed the difference between the NDC target and actual domestic emissions. \\
  \textbf{No ITMOs from unambitious NDCs} & Nat, BAU & ITMOs should not be bought from a seller with unambitious NDC target, i.e. a target above the BAU or both above cost-optimal and equal per capita 2\textdegree{C} emissions. \\
  \textbf{No ITMOs from %
  large emitter} & Nat & ITMOs should not be bought from a seller with %
  emissions above the world average. \\
    \textbf{Ambitious NDC in exchange for adequate climate finance} &  & Potential sellers would set their (conditional) NDC target below the BAU in exchange for adequate climate finance. \\
  \textbf{Conditional NDC prevailing in case of adequate climate finance} &  & The adequacy of a seller's NDC (to sell ITMOs) would be assessed using the conditional NDC target instead of the unconditional one iff the climate finance it receives matches the condition specified in its NDC. \\ %
\textbf{ITMO coalition with a joint NDC} &  & A coalition of countries would submit a joint NDC aligned with the Paris target and exchange ITMOs only among themselves. The NDC would ideally specify a carbon budget and yearly national targets. \\
  \bottomrule\\[-0.81em]
\end{tabular}}
\end{table}

\section{Conclusion\label{sec:conclusion}}

This paper reviewed and developed proposals for aligning Internationally Transferred Mitigation Outcomes (ITMOs) with the Paris temperature target. The sum of current NDC targets exceeds the required emissions level; ITMOs will therefore remain under-priced, and their unfettered use risks weakening overall ambition.

By restricting the use of ITMOs based on the adequacy of the seller's NDC, existing proposals (by e.g. \citealp{michaelowa_additionality_2019}) would limit hot air, but ITMOs would remain under-priced. To fully restore ambition and fairness, I proposed that a coalition of the willing submit a joint NDC aligned with the Paris target, and commit to exchanging ITMOs only among coalition members. I also suggested intermediary steps, such as strengthening reporting requirements to harmonize NDCs and establishing the principle that no ITMO should be sold by a country failing to meet its own NDC.

A tighter regulation of ITMOs combined with an ambitiously expanded JETP program \citep{steckel_climate_2017,bolton_why_2025} offers a viable pathway for effective international decarbonization. In both cases, a coalition of the willing seems the best option to reap the gains from the cooperation of ambitious countries. 

  %

\clearpage
\small
\renewcommand{\url}[1]{\href{#1}{Link}} %
\bibliographystyle{plainnaturl_clean} %
\bibliography{global_tax_attitudes}

\appendix

\renewcommand{\thetable}{A\arabic{table}}
\renewcommand{\thefigure}{A\arabic{figure}}
\setcounter{figure}{0}
\setcounter{table}{0}

\end{document}